% =========================================================
% Nested Interval Property
% =========================================================

\subsection*{Nested Interval Property}

\begin{lemmabox}{Lemma (Nested Closed Intervals Have Ordered Endpoints)}
\begin{lemma}[Nested Closed Intervals Have Ordered Endpoints]
\label{lem:nested-closed-intervals-have-ordered-endpoints}
Let $(a_n)$ and $(b_n)$ be sequences in $\mathbb{R}$ satisfying
\[
  (\forall n\in\mathbb{N})(a_n\leq b_n)
\]
and
\[
  (\forall n\in\mathbb{N})([a_{n+1},b_{n+1}]\subseteq [a_n,b_n]).
\]
Then, for all $m,n\in\mathbb{N}$,
\[
  n\leq m
  \Longrightarrow
  a_n\leq a_m\leq b_m\leq b_n.
\]
\hyperref[prf:nested-closed-intervals-have-ordered-endpoints]{Go to proof.}
\end{lemma}
\end{lemmabox}

\begin{remark*}[Standard quantified statement]
\[
\forall n\in\mathbb{N}\;
\Bigl([a_{n+1},b_{n+1}]\subseteq [a_n,b_n]
\Longrightarrow
a_n\le a_{n+1}\le b_{n+1}\le b_n\Bigr).
\]
\end{remark*}

\begin{remark*}[Predicate reading]
\[
[a_{n+1},b_{n+1}]\subseteq [a_n,b_n]
\Longrightarrow
a_n\le a_{n+1}\le b_{n+1}\le b_n.
\]
\end{remark*}

\begin{remark*}[Negated quantified statement]
\[
\exists n\in\mathbb{N}\;
\Bigl([a_{n+1},b_{n+1}]\subseteq [a_n,b_n]\land
(a_{n+1}<a_n\lor b_n<b_{n+1})\Bigr).
\]
\end{remark*}

\begin{remark*}[Negation predicate reading]
\[
(a_{n+1}<a_n\lor b_n<b_{n+1})
\Longrightarrow
[a_{n+1},b_{n+1}]\not\subseteq [a_n,b_n].
\]
\end{remark*}

\begin{remark*}[Failure modes]
\begin{description}
  \item[Exposition.]
  Failure is witnessed by the displayed negated or contrapositive condition.

  \item[\exists-condition.]
  \[
\exists n\in\mathbb{N}\;
\Bigl([a_{n+1},b_{n+1}]\subseteq [a_n,b_n]\land
(a_{n+1}<a_n\lor b_n<b_{n+1})\Bigr).
\]
\[
(a_{n+1}<a_n\lor b_n<b_{n+1})
\Longrightarrow
[a_{n+1},b_{n+1}]\not\subseteq [a_n,b_n].
\]
\end{description}
\end{remark*}

\begin{remark*}[Interpretation]
One-step nesting of closed intervals moves the left endpoint to the right
and the right endpoint to the left.
\end{remark*}

\begin{dependencies}
\hyperref[def:closed-interval]{Closed Interval}.
\end{dependencies}

\begin{lemmabox}{Lemma (Left Endpoints of Nested Intervals Are Increasing)}
\begin{lemma}[Left Endpoints of Nested Intervals Are Increasing]
\label{lem:left-endpoints-nested-intervals-increasing}
Let $(a_n)$ and $(b_n)$ be sequences in $\mathbb{R}$ such that
\[
  (\forall n\in\mathbb{N})(a_n\leq b_n)
\]
and
\[
  (\forall n\in\mathbb{N})([a_{n+1},b_{n+1}]\subseteq [a_n,b_n]).
\]
Then
\[
  (\forall n\in\mathbb{N})(a_n\leq a_{n+1}).
\]
\hyperref[prf:left-endpoints-nested-intervals-increasing]{Go to proof.}
\end{lemma}
\end{lemmabox}

\begin{remark*}[Standard quantified statement]
\[
\forall n\in\mathbb{N}\;(a_n\le a_{n+1}).
\]
\end{remark*}

\begin{remark*}[Predicate reading]
\[
F=([a_n,b_n])_{n\in\mathbb{N}},\qquad
a\colon\mathbb{N}\longrightarrow\mathbb{R},\qquad a(n)=a_n,
\]
\[
P_{\mathbb{N}}=\mathsf{OrderedSet}(\mathbb{N},\leq),\qquad
P_{\mathbb{R}}=\mathsf{OrderedSet}(\mathbb{R},\leq),
\]
\[
\operatorname{NestedClosedIntervals}(F,P_{\mathbb{R}})
\Longrightarrow
\operatorname{OrderPreserving}(a,P_{\mathbb{N}},P_{\mathbb{R}}).
\]
\end{remark*}

\begin{remark*}[Negated quantified statement]
\[
\exists n\in\mathbb{N}\;(a_{n+1}<a_n).
\]
\end{remark*}

\begin{remark*}[Negation predicate reading]
\[
F=([a_n,b_n])_{n\in\mathbb{N}},\qquad
a\colon\mathbb{N}\longrightarrow\mathbb{R},\qquad a(n)=a_n,
\]
\[
P_{\mathbb{N}}=\mathsf{OrderedSet}(\mathbb{N},\leq),\qquad
P_{\mathbb{R}}=\mathsf{OrderedSet}(\mathbb{R},\leq),
\]
\[
\neg\operatorname{OrderPreserving}(a,P_{\mathbb{N}},P_{\mathbb{R}})
\Longrightarrow
\neg\operatorname{NestedClosedIntervals}(F,P_{\mathbb{R}}).
\]
\end{remark*}

\begin{remark*}[Failure modes]
\begin{description}
  \item[Exposition.]
  Failure is witnessed by the displayed negated or contrapositive condition.

  \item[\exists-condition.]
  \[
\exists n\in\mathbb{N}\;(a_{n+1}<a_n).
\]
\[
F=([a_n,b_n])_{n\in\mathbb{N}},\qquad
a\colon\mathbb{N}\longrightarrow\mathbb{R},\qquad a(n)=a_n,
\]
\end{description}
\end{remark*}

\begin{remark*}[Interpretation]
In a nested family, the left endpoints cannot move left as the intervals
shrink.
\end{remark*}

\begin{dependencies}
\hyperref[def:closed-interval]{Closed Interval}.
\end{dependencies}

\begin{lemmabox}{Lemma (Right Endpoints of Nested Intervals Are Decreasing)}
\begin{lemma}[Right Endpoints of Nested Intervals Are Decreasing]
\label{lem:right-endpoints-nested-intervals-decreasing}
Let $(a_n)$ and $(b_n)$ be sequences in $\mathbb{R}$ such that
\[
  (\forall n\in\mathbb{N})(a_n\leq b_n)
\]
and
\[
  (\forall n\in\mathbb{N})([a_{n+1},b_{n+1}]\subseteq [a_n,b_n]).
\]
Then
\[
  (\forall n\in\mathbb{N})(b_{n+1}\leq b_n).
\]
\hyperref[prf:right-endpoints-nested-intervals-decreasing]{Go to proof.}
\end{lemma}
\end{lemmabox}

\begin{remark*}[Standard quantified statement]
\[
\forall n\in\mathbb{N}\;(b_{n+1}\le b_n).
\]
\end{remark*}

\begin{remark*}[Predicate reading]
\[
F=([a_n,b_n])_{n\in\mathbb{N}},\qquad
b\colon\mathbb{N}\longrightarrow\mathbb{R},\qquad b(n)=b_n,
\]
\[
P_{\mathbb{N}}=\mathsf{OrderedSet}(\mathbb{N},\leq),\qquad
P_{\mathbb{R}}=\mathsf{OrderedSet}(\mathbb{R},\leq),
\]
\[
\operatorname{NestedClosedIntervals}(F,P_{\mathbb{R}})
\Longrightarrow
\operatorname{OrderReversing}(b,P_{\mathbb{N}},P_{\mathbb{R}}).
\]
\end{remark*}

\begin{remark*}[Negated quantified statement]
\[
\exists n\in\mathbb{N}\;(b_n<b_{n+1}).
\]
\end{remark*}

\begin{remark*}[Negation predicate reading]
\[
F=([a_n,b_n])_{n\in\mathbb{N}},\qquad
b\colon\mathbb{N}\longrightarrow\mathbb{R},\qquad b(n)=b_n,
\]
\[
P_{\mathbb{N}}=\mathsf{OrderedSet}(\mathbb{N},\leq),\qquad
P_{\mathbb{R}}=\mathsf{OrderedSet}(\mathbb{R},\leq),
\]
\[
\neg\operatorname{OrderReversing}(b,P_{\mathbb{N}},P_{\mathbb{R}})
\Longrightarrow
\neg\operatorname{NestedClosedIntervals}(F,P_{\mathbb{R}}).
\]
\end{remark*}

\begin{remark*}[Failure modes]
\begin{description}
  \item[Exposition.]
  Failure is witnessed by the displayed negated or contrapositive condition.

  \item[\exists-condition.]
  \[
\exists n\in\mathbb{N}\;(b_n<b_{n+1}).
\]
\[
F=([a_n,b_n])_{n\in\mathbb{N}},\qquad
b\colon\mathbb{N}\longrightarrow\mathbb{R},\qquad b(n)=b_n,
\]
\end{description}
\end{remark*}

\begin{remark*}[Interpretation]
In a nested family, the right endpoints cannot move right as the intervals
shrink.
\end{remark*}

\begin{dependencies}
\hyperref[def:closed-interval]{Closed Interval}.
\end{dependencies}

\begin{lemmabox}{Lemma (Left Endpoints Are Bounded Above by Right Endpoints)}
\begin{lemma}[Left Endpoints Are Bounded Above by Right Endpoints]
\label{lem:left-endpoints-bounded-above-by-right-endpoints}
Let $(a_n)$ and $(b_n)$ be sequences in $\mathbb{R}$ such that
\[
  (\forall n\in\mathbb{N})(a_n\leq b_n)
\]
and
\[
  (\forall n\in\mathbb{N})([a_{n+1},b_{n+1}]\subseteq [a_n,b_n]).
\]
Then
\[
  (\forall m,n\in\mathbb{N})(a_n\leq b_m).
\]
Consequently, every $b_m$ is an upper bound of $\{a_n:n\in\mathbb{N}\}$, and
every $a_n$ is a lower bound of $\{b_m:m\in\mathbb{N}\}$.
\hyperref[prf:left-endpoints-bounded-above-by-right-endpoints]{Go to proof.}
\end{lemma}
\end{lemmabox}

\begin{remark*}[Standard quantified statement]
\[
\forall m,n\in\mathbb{N}\;(a_n\le b_m).
\]
\end{remark*}

\begin{remark*}[Predicate reading]
\[
F=([a_n,b_n])_{n\in\mathbb{N}},\qquad P_{\mathbb{R}}=\mathsf{OrderedSet}(\mathbb{R},\leq),
\]
\[
\operatorname{NestedClosedIntervals}(F,P_{\mathbb{R}})
\Longrightarrow
\forall m,n\in\mathbb{N}\;(a_n\le b_m).
\]
\end{remark*}

\begin{remark*}[Negated quantified statement]
\[
\exists m,n\in\mathbb{N}\;(b_m<a_n).
\]
\end{remark*}

\begin{remark*}[Negation predicate reading]
\[
F=([a_n,b_n])_{n\in\mathbb{N}},\qquad P_{\mathbb{R}}=\mathsf{OrderedSet}(\mathbb{R},\leq),
\]
\[
\exists m,n\in\mathbb{N}\;(b_m<a_n)
\Longrightarrow
\neg\operatorname{NestedClosedIntervals}(F,P_{\mathbb{R}}).
\]
\end{remark*}

\begin{remark*}[Failure modes]
\begin{description}
  \item[Exposition.]
  Failure is witnessed by the displayed negated or contrapositive condition.

  \item[\mathsf{OrderedSet}.]
  \[
\exists m,n\in\mathbb{N}\;(b_m<a_n).
\]
\[
F=([a_n,b_n])_{n\in\mathbb{N}},\qquad P_{\mathbb{R}}=\mathsf{OrderedSet}(\mathbb{R},\leq),
\]
\end{description}
\end{remark*}

\begin{remark*}[Interpretation]
Every left endpoint lies below every right endpoint, even when the two
endpoints come from different stages of the nested family.
\end{remark*}

\begin{dependencies}
\hyperref[def:closed-interval]{Closed Interval}.
\end{dependencies}

\begin{lemmabox}{Lemma (Endpoint Supremum Infimum Inequality)}
\begin{lemma}[Endpoint Supremum Infimum Inequality]
\label{lem:endpoint-supremum-infimum-inequality}
Let $(a_n)$ and $(b_n)$ be sequences in $\mathbb{R}$ satisfying
\[
  (\forall n\in\mathbb{N})(a_n\leq b_n)
\]
and
\[
  (\forall n\in\mathbb{N})([a_{n+1},b_{n+1}]\subseteq [a_n,b_n]).
\]
Define
\[
  \alpha:=\sup\{a_n:n\in\mathbb{N}\}
\]
and
\[
  \beta:=\inf\{b_n:n\in\mathbb{N}\}.
\]
Then
\[
  \alpha\leq\beta.
\]
\hyperref[prf:endpoint-supremum-infimum-inequality]{Go to proof.}
\end{lemma}
\end{lemmabox}

\begin{remark*}[Standard quantified statement]
\[
\alpha=\sup\{a_n:n\in\mathbb{N}\}\land
\beta=\inf\{b_n:n\in\mathbb{N}\}
\Longrightarrow
\alpha\le \beta.
\]
\end{remark*}

\begin{remark*}[Predicate reading]
\[
F=([a_n,b_n])_{n\in\mathbb{N}},\qquad
A=\{a_n:n\in\mathbb{N}\},\qquad B=\{b_n:n\in\mathbb{N}\},
\]
\[
P_{\mathbb{R}}=\mathsf{OrderedSet}(\mathbb{R},\leq),\qquad
\operatorname{NestedClosedIntervals}(F,P_{\mathbb{R}})
\land
\operatorname{LeastUpperBound}(\alpha,A,P_{\mathbb{R}})
\land
\operatorname{GreatestLowerBound}(\beta,B,P_{\mathbb{R}})
\Longrightarrow
\alpha\le \beta.
\]
\end{remark*}

\begin{remark*}[Negated quantified statement]
\[
\operatorname{NestedClosedIntervals}([a_n,b_n])
\land
\alpha=\sup\{a_n:n\in\mathbb{N}\}
\land
\beta=\inf\{b_n:n\in\mathbb{N}\}
\land
\beta<\alpha.
\]
\end{remark*}

\begin{remark*}[Negation predicate reading]
\[
F=([a_n,b_n])_{n\in\mathbb{N}},\qquad
A=\{a_n:n\in\mathbb{N}\},\qquad B=\{b_n:n\in\mathbb{N}\},
\]
\[
P_{\mathbb{R}}=\mathsf{OrderedSet}(\mathbb{R},\leq),\qquad
\beta<\alpha
\Longrightarrow
\neg\operatorname{NestedClosedIntervals}(F,P_{\mathbb{R}})
\lor
\neg\operatorname{LeastUpperBound}(\alpha,A,P_{\mathbb{R}})
\lor
\neg\operatorname{GreatestLowerBound}(\beta,B,P_{\mathbb{R}}).
\]
\end{remark*}

\begin{remark*}[Failure modes]
\begin{description}
  \item[Exposition.]
  Failure is witnessed by the displayed negated or contrapositive condition.

  \item[\operatorname{NestedClosedIntervals}.]
  \[
\operatorname{NestedClosedIntervals}([a_n,b_n])
\land
\alpha=\sup\{a_n:n\in\mathbb{N}\}
\land
\beta=\inf\{b_n:n\in\mathbb{N}\}
\land
\beta<\alpha.
\]
\[
F=([a_n,b_n])_{n\in\mathbb{N}},\qquad
A=\{a_n:n\in\mathbb{N}\},\qquad B=\{b_n:n\in\mathbb{N}\},
\]
\end{description}
\end{remark*}

\begin{remark*}[Interpretation]
The supremum of all left endpoints cannot pass the infimum of all right
endpoints.
\end{remark*}

\begin{dependencies}
\hyperref[def:closed-interval]{Closed Interval};
\hyperref[def:supremum]{Supremum};
\hyperref[def:infimum]{Infimum}.
\end{dependencies}

\begin{figure}[H]
\centering
\begin{tikzpicture}[x=1.15cm, y=0.8cm, >=Latex]
  \draw[->] (-0.4,0) -- (7.2,0) node[right] {$\mathbb R$};
  \foreach \x/\lab in {0.6/$a_1$,1.4/$a_2$,2.1/$a_3$,3.3/$\alpha$,4.7/$\beta$,5.6/$b_3$,6.3/$b_2$,6.9/$b_1$} {
    \draw (\x,0.08) -- (\x,-0.08) node[below=3pt] {\lab};
  }
  \foreach \y/\l/\r/\name in {1.8/0.6/6.9/$I_1$,1.25/1.4/6.3/$I_2$,0.7/2.1/5.6/$I_3$} {
    \draw[line width=1.1pt] (\l,\y) -- (\r,\y);
    \fill (\l,\y) circle (1.4pt);
    \fill (\r,\y) circle (1.4pt);
    \node[left] at (\l,\y) {\name};
  }
  \draw[dashed] (3.3,-0.05) -- (3.3,2.1);
  \draw[dashed] (4.7,-0.05) -- (4.7,2.1);
  \draw[decorate,decoration={brace,amplitude=4pt}] (3.3,2.25) -- (4.7,2.25)
    node[midway,above=5pt] {$\sup\{a_n\}\le \inf\{b_n\}$};
\end{tikzpicture}
\caption{Nested endpoints force every left endpoint below every right endpoint, so the supremum of the left endpoints lies no farther right than the infimum of the right endpoints.}
\label{fig:nested-interval-endpoints}
\end{figure}


\begin{theorembox}{Theorem (Nested Interval Theorem)}
\begin{theorem}[Nested Interval Theorem]
\label{thm:nested-interval-theorem}
Let $(a_n)$ and $(b_n)$ be sequences in $\mathbb{R}$ such that
\[
  (\forall n\in\mathbb{N})(a_n\leq b_n)
\]
and
\[
  (\forall n\in\mathbb{N})([a_{n+1},b_{n+1}]\subseteq [a_n,b_n]).
\]
Define
\[
  \alpha:=\sup\{a_n:n\in\mathbb{N}\},
  \qquad
  \beta:=\inf\{b_n:n\in\mathbb{N}\}.
\]
Then
\[
  \alpha\leq\beta
\]
and
\[
  \bigcap_{n=1}^{\infty}[a_n,b_n]\neq\varnothing.
\]
In particular, there exists $x\in\mathbb{R}$ such that
\[
  (\forall n\in\mathbb{N})(a_n\leq x\leq b_n).
\]
\hyperref[prf:nested-interval-theorem]{Go to proof.}
\end{theorem}
\end{theorembox}

\begin{remark*}[Standard quantified statement]
\[
\begin{aligned}
&\forall (a_n)\subseteq\mathbb{R}\;\forall (b_n)\subseteq\mathbb{R}\;\\
&\Bigl[
\bigl(\forall n\in\mathbb{N}\;(a_n\le b_n)\bigr)
\land
\bigl(\forall n\in\mathbb{N}\;([a_{n+1},b_{n+1}]\subseteq [a_n,b_n])\bigr)
\\
&\qquad\Longrightarrow
\exists x\in\mathbb{R}\;\forall n\in\mathbb{N}\;(x\in [a_n,b_n])
\Bigr].
\end{aligned}
\]
\end{remark*}

\begin{remark*}[Predicate reading]
\[
F=(I_n)_{n\in\mathbb{N}},\qquad C=\bigcap_{n=1}^{\infty}I_n,\qquad
P_{\mathbb{R}}=\mathsf{OrderedSet}(\mathbb{R},\leq),
\]
\[
\operatorname{NestedClosedBoundedIntervals}(F,P_{\mathbb{R}})\Longrightarrow
\operatorname{IsNonempty}(C).
\]
\end{remark*}

\begin{remark*}[Negated quantified statement]
\[
\begin{aligned}
&\exists (a_n)\subseteq\mathbb{R}\;\exists (b_n)\subseteq\mathbb{R}\;\\
&\Bigl[
\bigl(\forall n\in\mathbb{N}\;(a_n\le b_n)\bigr)
\land
\bigl(\forall n\in\mathbb{N}\;([a_{n+1},b_{n+1}]\subseteq [a_n,b_n])\bigr)
\land
\forall x\in\mathbb{R}\;\exists n\in\mathbb{N}\;(x\notin [a_n,b_n])
\Bigr].
\end{aligned}
\]
\end{remark*}

\begin{remark*}[Negation predicate reading]
\[
F=([a_n,b_n])_{n\in\mathbb{N}},\qquad C=\bigcap_{n=1}^{\infty}[a_n,b_n],\qquad
P_{\mathbb{R}}=\mathsf{OrderedSet}(\mathbb{R},\leq),
\]
\[
C=\varnothing
\Longrightarrow
\neg\operatorname{NestedClosedBoundedIntervals}(F,P_{\mathbb{R}}).
\]
\end{remark*}

\begin{remark*}[Failure modes]
\begin{description}
  \item[Exposition.]
  Failure is witnessed by the displayed negated or contrapositive condition.

  \item[\mathsf{OrderedSet}.]
  \[
\begin{aligned}
&\exists (a_n)\subseteq\mathbb{R}\;\exists (b_n)\subseteq\mathbb{R}\;\\
&\Bigl[
\bigl(\forall n\in\mathbb{N}\;(a_n\le b_n)\bigr)
\land
\bigl(\forall n\in\mathbb{N}\;([a_{n+1},b_{n+1}]\subseteq [a_n,b_n])\bigr)
\land
\forall x\in\mathbb{R}\;\exists n\in\mathbb{N}\;(x\notin [a_n,b_n])
\Bigr].
\end{aligned}
\]
\[
F=([a_n,b_n])_{n\in\mathbb{N}},\qquad C=\bigcap_{n=1}^{\infty}[a_n,b_n],\qquad
P_{\mathbb{R}}=\mathsf{OrderedSet}(\mathbb{R},\leq),
\]
\end{description}
\end{remark*}

\begin{remark*}[Interpretation]
Completeness prevents a nested closed bounded family from losing every point
at the limiting stage.
\end{remark*}

\begin{dependencies}
\hyperref[def:closed-interval]{Closed Interval};
\hyperref[def:bounded-subset-of-r]{Bounded Subset of the Real Line};
\hyperref[def:supremum]{Supremum};
\hyperref[ax:completeness-of-reals]{Axiom~\ref*{ax:completeness-of-reals}}.
\end{dependencies}

\begin{corollarybox}{Corollary (Nested Interval Intersection Contains the Endpoint Supremum)}
\begin{corollary}[Nested Interval Intersection Contains the Endpoint Supremum]
\label{cor:nested-interval-intersection-contains-endpoint-supremum}
Let $(a_n)$ and $(b_n)$ satisfy the hypotheses of the Nested Interval Theorem,
and define
\[
  \alpha:=\sup\{a_n:n\in\mathbb{N}\}.
\]
Then
\[
  (\forall n\in\mathbb{N})(a_n\leq\alpha\leq b_n).
\]
Therefore,
\[
  \alpha\in \bigcap_{n=1}^{\infty}[a_n,b_n].
\]
\hyperref[prf:nested-interval-intersection-contains-endpoint-supremum]{Go to proof.}
\end{corollary}
\end{corollarybox}

\begin{remark*}[Standard quantified statement]
\[
\alpha=\sup\{a_n:n\in\mathbb{N}\}
\Longrightarrow
\forall n\in\mathbb{N}\;(\alpha\in [a_n,b_n]).
\]
\end{remark*}

\begin{remark*}[Predicate reading]
\[
F=([a_n,b_n])_{n\in\mathbb{N}},\qquad A=\{a_n:n\in\mathbb{N}\},\qquad
C=\bigcap_{n=1}^{\infty}[a_n,b_n],
\]
\[
P_{\mathbb{R}}=\mathsf{OrderedSet}(\mathbb{R},\leq),\qquad
\operatorname{NestedClosedBoundedIntervals}(F,P_{\mathbb{R}})
\land
\operatorname{LeastUpperBound}(\alpha,A,P_{\mathbb{R}})
\Longrightarrow
\alpha\in C.
\]
\end{remark*}

\begin{remark*}[Negated quantified statement]
\[
\operatorname{NestedClosedBoundedIntervals}([a_n,b_n])
\land
\alpha=\sup\{a_n:n\in\mathbb{N}\}
\land
\exists n\in\mathbb{N}\;(\alpha\notin [a_n,b_n]).
\]
\end{remark*}

\begin{remark*}[Negation predicate reading]
\[
F=([a_n,b_n])_{n\in\mathbb{N}},\qquad A=\{a_n:n\in\mathbb{N}\},\qquad
C=\bigcap_{n=1}^{\infty}[a_n,b_n],
\]
\[
P_{\mathbb{R}}=\mathsf{OrderedSet}(\mathbb{R},\leq),\qquad
\alpha\notin C
\Longrightarrow
\neg\operatorname{LeastUpperBound}(\alpha,A,P_{\mathbb{R}})
\lor
\neg\operatorname{NestedClosedBoundedIntervals}(F,P_{\mathbb{R}}).
\]
\end{remark*}

\begin{remark*}[Failure modes]
\begin{description}
  \item[Exposition.]
  Failure is witnessed by the displayed negated or contrapositive condition.

  \item[\operatorname{NestedClosedBoundedIntervals}.]
  \[
\operatorname{NestedClosedBoundedIntervals}([a_n,b_n])
\land
\alpha=\sup\{a_n:n\in\mathbb{N}\}
\land
\exists n\in\mathbb{N}\;(\alpha\notin [a_n,b_n]).
\]
\[
F=([a_n,b_n])_{n\in\mathbb{N}},\qquad A=\{a_n:n\in\mathbb{N}\},\qquad
C=\bigcap_{n=1}^{\infty}[a_n,b_n],
\]
\end{description}
\end{remark*}

\begin{remark*}[Interpretation]
The least upper bound of the left endpoints is itself retained by every
interval in the nested family.
\end{remark*}

\begin{dependencies}
\hyperref[def:closed-interval]{Closed Interval};
\hyperref[def:bounded-subset-of-r]{Bounded Subset of the Real Line};
\hyperref[def:supremum]{Supremum}.
\end{dependencies}

\begin{corollarybox}{Corollary (Nested Interval Intersection Contains the Endpoint Infimum)}
\begin{corollary}[Nested Interval Intersection Contains the Endpoint Infimum]
\label{cor:nested-interval-intersection-contains-endpoint-infimum}
Let $(a_n)$ and $(b_n)$ satisfy the hypotheses of the Nested Interval Theorem,
and define
\[
  \beta:=\inf\{b_n:n\in\mathbb{N}\}.
\]
Then
\[
  (\forall n\in\mathbb{N})(a_n\leq\beta\leq b_n).
\]
Therefore,
\[
  \beta\in \bigcap_{n=1}^{\infty}[a_n,b_n].
\]
\hyperref[prf:nested-interval-intersection-contains-endpoint-infimum]{Go to proof.}
\end{corollary}
\end{corollarybox}

\begin{remark*}[Standard quantified statement]
\[
\beta=\inf\{b_n:n\in\mathbb{N}\}
\Longrightarrow
\forall n\in\mathbb{N}\;(\beta\in [a_n,b_n]).
\]
\end{remark*}

\begin{remark*}[Predicate reading]
\[
F=([a_n,b_n])_{n\in\mathbb{N}},\qquad B=\{b_n:n\in\mathbb{N}\},\qquad
C=\bigcap_{n=1}^{\infty}[a_n,b_n],
\]
\[
P_{\mathbb{R}}=\mathsf{OrderedSet}(\mathbb{R},\leq),\qquad
\operatorname{NestedClosedBoundedIntervals}(F,P_{\mathbb{R}})
\land
\operatorname{GreatestLowerBound}(\beta,B,P_{\mathbb{R}})
\Longrightarrow
\beta\in C.
\]
\end{remark*}

\begin{remark*}[Negated quantified statement]
\[
\operatorname{NestedClosedBoundedIntervals}([a_n,b_n])
\land
\beta=\inf\{b_n:n\in\mathbb{N}\}
\land
\exists n\in\mathbb{N}\;(\beta\notin [a_n,b_n]).
\]
\end{remark*}

\begin{remark*}[Negation predicate reading]
\[
F=([a_n,b_n])_{n\in\mathbb{N}},\qquad B=\{b_n:n\in\mathbb{N}\},\qquad
C=\bigcap_{n=1}^{\infty}[a_n,b_n],
\]
\[
P_{\mathbb{R}}=\mathsf{OrderedSet}(\mathbb{R},\leq),\qquad
\beta\notin C
\Longrightarrow
\neg\operatorname{GreatestLowerBound}(\beta,B,P_{\mathbb{R}})
\lor
\neg\operatorname{NestedClosedBoundedIntervals}(F,P_{\mathbb{R}}).
\]
\end{remark*}

\begin{remark*}[Failure modes]
\begin{description}
  \item[Exposition.]
  Failure is witnessed by the displayed negated or contrapositive condition.

  \item[\operatorname{NestedClosedBoundedIntervals}.]
  \[
\operatorname{NestedClosedBoundedIntervals}([a_n,b_n])
\land
\beta=\inf\{b_n:n\in\mathbb{N}\}
\land
\exists n\in\mathbb{N}\;(\beta\notin [a_n,b_n]).
\]
\[
F=([a_n,b_n])_{n\in\mathbb{N}},\qquad B=\{b_n:n\in\mathbb{N}\},\qquad
C=\bigcap_{n=1}^{\infty}[a_n,b_n],
\]
\end{description}
\end{remark*}

\begin{remark*}[Interpretation]
The greatest lower bound of the right endpoints is itself retained by every
interval in the nested family.
\end{remark*}

\begin{dependencies}
\hyperref[def:closed-interval]{Closed Interval};
\hyperref[def:bounded-subset-of-r]{Bounded Subset of the Real Line};
\hyperref[def:infimum]{Infimum}.
\end{dependencies}

\begin{corollarybox}{Corollary (Nested Intervals with Vanishing Length Have Unique Point)}
\begin{corollary}[Nested Intervals with Vanishing Length Have Unique Point]
\label{cor:nested-intervals-vanishing-length-have-unique-point}
Let $(a_n)$ and $(b_n)$ be sequences in $\mathbb{R}$ such that
\[
  (\forall n\in\mathbb{N})(a_n\leq b_n),
\]
\[
  (\forall n\in\mathbb{N})([a_{n+1},b_{n+1}]\subseteq [a_n,b_n]),
\]
and
\[
  (\forall\varepsilon>0)(\exists N\in\mathbb{N})(\forall n\in\mathbb{N})
  \bigl(n\geq N\Longrightarrow 0\leq b_n-a_n<\varepsilon\bigr).
\]
Then there exists a unique $x\in\mathbb{R}$ such that
\[
  (\forall n\in\mathbb{N})(a_n\leq x\leq b_n).
\]
Equivalently,
\[
  \bigcap_{n=1}^{\infty}[a_n,b_n]=\{x\}.
\]
Moreover,
\[
  x=\sup\{a_n:n\in\mathbb{N}\}=\inf\{b_n:n\in\mathbb{N}\}.
\]
\hyperref[prf:nested-intervals-vanishing-length-have-unique-point]{Go to proof.}
\end{corollary}
\end{corollarybox}

\begin{remark*}[Standard quantified statement]
\[
\lim_{n\longrightarrow\infty}(b_n-a_n)=0
\Longrightarrow
\exists x\in\mathbb{R}\;\forall y\in\mathbb{R}\;
\left(y\in\bigcap_{n=1}^{\infty}[a_n,b_n]\Longleftrightarrow y=x\right).
\]
\end{remark*}

\begin{remark*}[Predicate reading]
\[
F=([a_n,b_n])_{n\in\mathbb{N}},\qquad
d\colon\mathbb{N}\longrightarrow\mathbb{R},\qquad d_n=b_n-a_n,
\]
\[
P_{\mathbb{R}}=\mathsf{OrderedSet}(\mathbb{R},\leq),\qquad X=\mathbb{R},\qquad
\operatorname{NestedClosedBoundedIntervals}(F,P_{\mathbb{R}})
\land
\operatorname{ConvergesTo}(d,0,X)
\Longrightarrow
\exists x\in\mathbb{R}\;
\left(\bigcap_{n=1}^{\infty}[a_n,b_n]=\{x\}\right).
\]
\end{remark*}

\begin{remark*}[Negated quantified statement]
\[
\operatorname{NestedClosedBoundedIntervals}([a_n,b_n])
\land
\lim_{n\longrightarrow\infty}(b_n-a_n)=0
\land
\forall x\in\mathbb{R}\;\exists y\in\mathbb{R}\;
\left(y\in\bigcap_{n=1}^{\infty}[a_n,b_n]\land y\ne x\right).
\]
\end{remark*}

\begin{remark*}[Negation predicate reading]
\[
F=([a_n,b_n])_{n\in\mathbb{N}},\qquad
d\colon\mathbb{N}\longrightarrow\mathbb{R},\qquad d_n=b_n-a_n,
\]
\[
P_{\mathbb{R}}=\mathsf{OrderedSet}(\mathbb{R},\leq),\qquad X=\mathbb{R},\qquad
\neg\exists x\in\mathbb{R}\;
\left(\bigcap_{n=1}^{\infty}[a_n,b_n]=\{x\}\right)
\Longrightarrow
\neg\operatorname{NestedClosedBoundedIntervals}(F,P_{\mathbb{R}})
\lor
\neg\operatorname{ConvergesTo}(d,0,X).
\]
\end{remark*}

\begin{remark*}[Failure modes]
\begin{description}
  \item[Exposition.]
  Failure is witnessed by the displayed negated or contrapositive condition.

  \item[\operatorname{NestedClosedBoundedIntervals}.]
  \[
\operatorname{NestedClosedBoundedIntervals}([a_n,b_n])
\land
\lim_{n\longrightarrow\infty}(b_n-a_n)=0
\land
\forall x\in\mathbb{R}\;\exists y\in\mathbb{R}\;
\left(y\in\bigcap_{n=1}^{\infty}[a_n,b_n]\land y\ne x\right).
\]
\[
F=([a_n,b_n])_{n\in\mathbb{N}},\qquad
d\colon\mathbb{N}\longrightarrow\mathbb{R},\qquad d_n=b_n-a_n,
\]
\end{description}
\end{remark*}

\begin{remark*}[Interpretation]
When nested lengths tend to zero, the surviving intersection cannot contain
two distinct points.
\end{remark*}

\begin{dependencies}
\hyperref[def:closed-interval]{Closed Interval};
\hyperref[def:bounded-subset-of-r]{Bounded Subset of the Real Line}.
\end{dependencies}

\begin{corollarybox}{Corollary (Nested Intervals with Vanishing Length Have Equal Endpoint Limits)}
\begin{corollary}[Nested Intervals with Vanishing Length Have Equal Endpoint Limits]
\label{cor:nested-intervals-vanishing-length-have-equal-endpoint-limits}
Let $(a_n)$ and $(b_n)$ satisfy
\[
  (\forall n\in\mathbb{N})(a_n\leq b_n),
\]
\[
  (\forall n\in\mathbb{N})([a_{n+1},b_{n+1}]\subseteq [a_n,b_n]),
\]
and
\[
  (\forall\varepsilon>0)(\exists N\in\mathbb{N})(\forall n\in\mathbb{N})
  \bigl(n\geq N\Longrightarrow b_n-a_n<\varepsilon\bigr).
\]
Then there exists $x\in\mathbb{R}$ such that
\[
  \lim_{n\longrightarrow\infty}a_n=x
  \qquad\text{and}\qquad
  \lim_{n\longrightarrow\infty}b_n=x.
\]
\hyperref[prf:nested-intervals-vanishing-length-have-equal-endpoint-limits]{Go to proof.}
\end{corollary}
\end{corollarybox}

\begin{remark*}[Standard quantified statement]
\[
\lim_{n\longrightarrow\infty}(b_n-a_n)=0
\Longrightarrow
\exists L\in\mathbb{R}\;
\left(\lim_{n\longrightarrow\infty}a_n=L\land \lim_{n\longrightarrow\infty}b_n=L\right).
\]
\end{remark*}

\begin{remark*}[Predicate reading]
\[
F=([a_n,b_n])_{n\in\mathbb{N}},\qquad
a\colon\mathbb{N}\longrightarrow\mathbb{R},\qquad b\colon\mathbb{N}\longrightarrow\mathbb{R},\qquad d\colon\mathbb{N}\longrightarrow\mathbb{R},
\]
\[
a(n)=a_n,\qquad b(n)=b_n,\qquad d_n=b_n-a_n,
\]
\[
P_{\mathbb{R}}=\mathsf{OrderedSet}(\mathbb{R},\leq),\qquad X=\mathbb{R},\qquad
\operatorname{NestedClosedBoundedIntervals}(F,P_{\mathbb{R}})
\land
\operatorname{ConvergesTo}(d,0,X)
\Longrightarrow
\exists L\in\mathbb{R}\;
\left(
\operatorname{ConvergesTo}(a,L,X)
\land
\operatorname{ConvergesTo}(b,L,X)
\right).
\]
\end{remark*}

\begin{remark*}[Negated quantified statement]
\[
\operatorname{NestedClosedBoundedIntervals}([a_n,b_n])
\land
\lim_{n\longrightarrow\infty}(b_n-a_n)=0
\land
\forall L\in\mathbb{R}\;
\left(\lim_{n\longrightarrow\infty}a_n\ne L\lor \lim_{n\longrightarrow\infty}b_n\ne L\right).
\]
\end{remark*}

\begin{remark*}[Negation predicate reading]
\[
a\colon\mathbb{N}\longrightarrow\mathbb{R},\qquad b\colon\mathbb{N}\longrightarrow\mathbb{R},\qquad
d\colon\mathbb{N}\longrightarrow\mathbb{R},\qquad d_n=b_n-a_n,
\]
\[
F=([a_n,b_n])_{n\in\mathbb{N}},\qquad
P_{\mathbb{R}}=\mathsf{OrderedSet}(\mathbb{R},\leq),\qquad X=\mathbb{R},\qquad
\neg\exists L\in\mathbb{R}\;
\left(
\operatorname{ConvergesTo}(a,L,X)
\land
\operatorname{ConvergesTo}(b,L,X)
\right)
\Longrightarrow
\neg\operatorname{NestedClosedBoundedIntervals}(F,P_{\mathbb{R}})
\lor
\neg\operatorname{ConvergesTo}(d,0,X).
\]
\end{remark*}

\begin{remark*}[Failure modes]
\begin{description}
  \item[Exposition.]
  Failure is witnessed by the displayed negated or contrapositive condition.

  \item[\operatorname{NestedClosedBoundedIntervals}.]
  \[
\operatorname{NestedClosedBoundedIntervals}([a_n,b_n])
\land
\lim_{n\longrightarrow\infty}(b_n-a_n)=0
\land
\forall L\in\mathbb{R}\;
\left(\lim_{n\longrightarrow\infty}a_n\ne L\lor \lim_{n\longrightarrow\infty}b_n\ne L\right).
\]
\[
a\colon\mathbb{N}\longrightarrow\mathbb{R},\qquad b\colon\mathbb{N}\longrightarrow\mathbb{R},\qquad
d\colon\mathbb{N}\longrightarrow\mathbb{R},\qquad d_n=b_n-a_n,
\]
\end{description}
\end{remark*}

\begin{remark*}[Interpretation]
Vanishing interval length forces the two endpoint sequences to approach the
same real number.
\end{remark*}

\begin{dependencies}
\hyperref[def:closed-interval]{Closed Interval};
\hyperref[def:bounded-subset-of-r]{Bounded Subset of the Real Line}.
\end{dependencies}

\begin{corollarybox}{Corollary (Unique Point in Nested Intervals Is Endpoint Limit)}
\begin{corollary}[Unique Point in Nested Intervals Is Endpoint Limit]
\label{cor:unique-point-in-nested-intervals-is-endpoint-limit}
Let $(a_n)$ and $(b_n)$ be sequences in $\mathbb{R}$ satisfying
\[
  (\forall n\in\mathbb{N})(a_n\leq b_n)
\]
and
\[
  (\forall n\in\mathbb{N})([a_{n+1},b_{n+1}]\subseteq [a_n,b_n]).
\]
Suppose there exists $x\in\mathbb{R}$ such that
\[
  \bigcap_{n=1}^{\infty}[a_n,b_n]=\{x\}.
\]
Then
\[
  \lim_{n\longrightarrow\infty}a_n=x
  \quad\text{and}\quad
  \lim_{n\longrightarrow\infty}b_n=x.
\]
\hyperref[prf:unique-point-in-nested-intervals-is-endpoint-limit]{Go to proof.}
\end{corollary}
\end{corollarybox}

\begin{remark*}[Standard quantified statement]
\[
\left[
\lim_{n\longrightarrow\infty}(b_n-a_n)=0
\land
\bigcap_{n=1}^{\infty}[a_n,b_n]=\{x\}
\right]
\Longrightarrow
\left(\lim_{n\longrightarrow\infty}a_n=x\land \lim_{n\longrightarrow\infty}b_n=x\right).
\]
\end{remark*}

\begin{remark*}[Predicate reading]
\[
a\colon\mathbb{N}\longrightarrow\mathbb{R},\qquad b\colon\mathbb{N}\longrightarrow\mathbb{R},\qquad d\colon\mathbb{N}\longrightarrow\mathbb{R},\qquad d_n=b_n-a_n,
\]
\[
C=\bigcap_{n=1}^{\infty}[a_n,b_n],\qquad X=\mathbb{R},\qquad
C=\{x\}
\land
\operatorname{ConvergesTo}(d,0,X)
\Longrightarrow
\operatorname{ConvergesTo}(a,x,X)
\land
\operatorname{ConvergesTo}(b,x,X).
\]
\end{remark*}

\begin{remark*}[Negated quantified statement]
\[
\bigcap_{n=1}^{\infty}[a_n,b_n]=\{x\}
\land
\lim_{n\longrightarrow\infty}(b_n-a_n)=0
\land
\left(\lim_{n\longrightarrow\infty}a_n\ne x\lor \lim_{n\longrightarrow\infty}b_n\ne x\right).
\]
\end{remark*}

\begin{remark*}[Negation predicate reading]
\[
a\colon\mathbb{N}\longrightarrow\mathbb{R},\qquad b\colon\mathbb{N}\longrightarrow\mathbb{R},\qquad d\colon\mathbb{N}\longrightarrow\mathbb{R},\qquad d_n=b_n-a_n,
\]
\[
C=\bigcap_{n=1}^{\infty}[a_n,b_n],\qquad X=\mathbb{R},\qquad
\left(\lim_{n\longrightarrow\infty}a_n\ne x\lor \lim_{n\longrightarrow\infty}b_n\ne x\right)
\Longrightarrow
C\ne\{x\}
\lor
\neg\operatorname{ConvergesTo}(d,0,X).
\]
\end{remark*}

\begin{remark*}[Failure modes]
\begin{description}
  \item[Exposition.]
  Failure is witnessed by the displayed negated or contrapositive condition.

  \item[\operatorname{IsEmpty}.]
  \[
\bigcap_{n=1}^{\infty}[a_n,b_n]=\{x\}
\land
\lim_{n\longrightarrow\infty}(b_n-a_n)=0
\land
\left(\lim_{n\longrightarrow\infty}a_n\ne x\lor \lim_{n\longrightarrow\infty}b_n\ne x\right).
\]
\[
a\colon\mathbb{N}\longrightarrow\mathbb{R},\qquad b\colon\mathbb{N}\longrightarrow\mathbb{R},\qquad d\colon\mathbb{N}\longrightarrow\mathbb{R},\qquad d_n=b_n-a_n,
\]
\end{description}
\end{remark*}

\begin{remark*}[Interpretation]
If the intersection has a single survivor, both endpoint sequences converge
to that survivor.
\end{remark*}

\begin{dependencies}
\hyperref[def:closed-interval]{Closed Interval};
\hyperref[def:bounded-subset-of-r]{Bounded Subset of the Real Line}.
\end{dependencies}

\begin{figure}[H]
\centering
\begin{tikzpicture}[x=1cm, y=0.8cm, >=Latex]
  \draw[->] (-0.2,1.7) -- (4.7,1.7) node[right] {$\mathbb R$};
  \node[left] at (-0.2,1.7) {open};
  \foreach \r/\y in {4.0/2.35,2.8/2.0,1.6/1.65,0.8/1.3} {
    \draw[line width=1pt] (0,\y) -- (\r,\y);
    \draw (0,\y) circle (1.4pt);
    \draw (\r,\y) circle (1.4pt);
  }
  \node[below] at (0,1.05) {$0$ omitted};
  \node[below] at (2.5,1.05) {$(0,1/n)$};

  \begin{scope}[yshift=-2.4cm]
    \draw[->] (-0.2,1.7) -- (6.2,1.7) node[right] {$\mathbb R$};
    \node[left] at (-0.2,1.7) {unbounded};
    \foreach \l/\y in {0.4/2.35,1.3/2.0,2.2/1.65,3.1/1.3} {
      \draw[line width=1pt] (\l,\y) -- (5.6,\y);
      \fill (\l,\y) circle (1.4pt);
      \draw[->,line width=1pt] (5.4,\y) -- (5.9,\y);
    }
    \node[below] at (3.0,1.05) {$[n,\infty)$ escapes right};
  \end{scope}
\end{tikzpicture}
\caption{Closedness and boundedness are separate hypotheses: open intervals can lose their endpoint, and unbounded closed intervals can drift away.}
\label{fig:nested-interval-failures}
\end{figure}


\begin{propositionbox}{Proposition (Open Nested Intervals Need Not Have Nonempty Intersection)}
\begin{proposition}[Open Nested Intervals Need Not Have Nonempty Intersection]
\label{prop:open-nested-intervals-need-not-have-nonempty-intersection}
There exists a sequence $(I_n)_{n\in\mathbb{N}}$ of nonempty bounded open
intervals in $\mathbb{R}$ such that
\[
  (\forall n\in\mathbb{N})(I_{n+1}\subseteq I_n),
\]
but
\[
  \bigcap_{n=1}^{\infty}I_n=\varnothing.
\]
One such sequence is
\[
  I_n=\left(0,\frac{1}{n}\right).
\]
\hyperref[prf:open-nested-intervals-need-not-have-nonempty-intersection]{Go to proof.}
\end{proposition}
\end{propositionbox}

\begin{remark*}[Standard quantified statement]
\[
\exists (a_n)\subseteq\mathbb{R}\;\exists (b_n)\subseteq\mathbb{R}\;
\left[
\begin{aligned}
&\forall n\in\mathbb{N}\;(a_n<b_n)
\land
\forall n\in\mathbb{N}\;((a_{n+1},b_{n+1})\subseteq (a_n,b_n))\\
&\land
\forall x\in\mathbb{R}\;\exists n\in\mathbb{N}\;(x\notin(a_n,b_n))
\end{aligned}
\right].
\]
\end{remark*}

\begin{remark*}[Predicate reading]
\[
\exists F\;\exists\mathcal{I}\;\exists C\;\bigl(
\mathcal{I}=\mathsf{IndexedFamily}(F,\mathbb{N},\mathbb{R})
\land C=\bigcap_{n=1}^{\infty}I_n
\land \Phi=\mathsf{OpenBoundedInterval}
\land
\operatorname{MembersSatisfy}(\mathcal{I},\Phi,\mathbb{R})
\land
\operatorname{Nested}(\mathcal{I},\subseteq)
\land
\operatorname{IsEmpty}(C)
\bigr).
\]
\end{remark*}

\begin{remark*}[Negated quantified statement]
\[
\forall (a_n),(b_n)\;\forall\mathcal{I}\;
\left[
\mathcal{I}=\mathsf{IndexedFamily}((a_n,b_n),\mathbb{N},\mathbb{R})
\land
\operatorname{MembersSatisfy}(\mathcal{I},\mathsf{OpenBoundedInterval},\mathbb{R})
\land
\operatorname{Nested}(\mathcal{I},\subseteq)
\Longrightarrow
\exists x\in\mathbb{R}\;\forall n\in\mathbb{N}\;(x\in(a_n,b_n))
\right].
\]
\end{remark*}

\begin{remark*}[Negation predicate reading]
\[
\forall (a_n),(b_n)\;\forall\mathcal{I}\;
\left[
\mathcal{I}=\mathsf{IndexedFamily}((a_n,b_n),\mathbb{N},\mathbb{R})
\land
\operatorname{MembersSatisfy}(\mathcal{I},\mathsf{OpenBoundedInterval},\mathbb{R})
\land
\operatorname{Nested}(\mathcal{I},\subseteq)
\Longrightarrow
\exists x\in\mathbb{R}\;\forall n\in\mathbb{N}\;(x\in(a_n,b_n))
\right].
\]

The negation predicate reading is the predicate-level form of the displayed negated statement.
\end{remark*}

\begin{remark*}[Failure modes]
\begin{description}
  \item[Exposition.]
  Failure is witnessed by the displayed negated or contrapositive condition.

  \item[\operatorname{MembersSatisfy} / \operatorname{Nested}.]
  \[
\forall (a_n),(b_n)\;\forall\mathcal{I}\;
\left[
\mathcal{I}=\mathsf{IndexedFamily}((a_n,b_n),\mathbb{N},\mathbb{R})
\land
\operatorname{MembersSatisfy}(\mathcal{I},\mathsf{OpenBoundedInterval},\mathbb{R})
\land
\operatorname{Nested}(\mathcal{I},\subseteq)
\Longrightarrow
\exists x\in\mathbb{R}\;\forall n\in\mathbb{N}\;(x\in(a_n,b_n))
\right].
\]
\[
\exists F\;\exists\mathcal{I}\;\exists C\;\bigl(
\mathcal{I}=\mathsf{IndexedFamily}(F,\mathbb{N},\mathbb{R})
\land C=\bigcap_{n=1}^{\infty}I_n
\land \Phi=\mathsf{OpenBoundedInterval}
\land
\operatorname{MembersSatisfy}(\mathcal{I},\Phi,\mathbb{R})
\land
\operatorname{Nested}(\mathcal{I},\subseteq)
\land
\operatorname{IsEmpty}(C)
\bigr).
\]
\end{description}
\end{remark*}

\begin{remark*}[Interpretation]
Open intervals can lose the limiting endpoint that every closed truncation
would retain.
\end{remark*}

\begin{dependencies}
\hyperref[def:open-interval]{Open Interval};
\hyperref[def:bounded-subset-of-r]{Bounded Subset of the Real Line}.
\end{dependencies}

\begin{propositionbox}{Proposition (Closedness Is Necessary in the Nested Interval Theorem)}
\begin{proposition}[Closedness Is Necessary in the Nested Interval Theorem]
\label{prop:closedness-is-necessary-in-nested-interval-theorem}
There exists a sequence $(I_n)_{n\in\mathbb{N}}$ of nonempty bounded intervals
in $\mathbb{R}$, not all of which are closed, such that
\[
  (\forall n\in\mathbb{N})(I_{n+1}\subseteq I_n),
\]
but
\[
  \bigcap_{n=1}^{\infty}I_n=\varnothing.
\]
For example,
\[
  I_n=\left(0,\frac{1}{n}\right].
\]
\hyperref[prf:closedness-is-necessary-in-nested-interval-theorem]{Go to proof.}
\end{proposition}
\end{propositionbox}

\begin{remark*}[Standard quantified statement]
\[
\exists I_n\;\exists\mathcal{I}\;
\bigl(
\mathcal{I}=\mathsf{IndexedFamily}(I_n,\mathbb{N},\mathbb{R})
\land
\operatorname{MembersSatisfy}(\mathcal{I},\mathsf{NonemptyBoundedInterval},\mathbb{R})
\land
\operatorname{Nested}(\mathcal{I},\subseteq)
\land
\neg\operatorname{IsClosedInterval}(I_n)
\land
\bigcap_{n=1}^{\infty}I_n=\varnothing
\bigr).
\]
\end{remark*}

\begin{remark*}[Predicate reading]
\[
\mathcal{I}=\mathsf{IndexedFamily}(I_n,\mathbb{N},\mathbb{R}),\qquad C=\bigcap_{n=1}^{\infty}I_n,\qquad
\neg\operatorname{IsClosedInterval}(I_n,\mathbb{R})
\quad\text{can make}\quad
\operatorname{IsNonempty}(C)
\text{ fail.}
\]
\end{remark*}

\begin{remark*}[Negated quantified statement]
\[
\forall I_n\;\forall\mathcal{I}\;
\left[
\mathcal{I}=\mathsf{IndexedFamily}(I_n,\mathbb{N},\mathbb{R})
\land
\operatorname{MembersSatisfy}(\mathcal{I},\mathsf{NonemptyBoundedInterval},\mathbb{R})
\land
\operatorname{Nested}(\mathcal{I},\subseteq)
\Longrightarrow
\bigcap_{n=1}^{\infty}I_n\ne\varnothing
\right].
\]
\end{remark*}

\begin{remark*}[Negation predicate reading]
\[
C=\bigcap_{n=1}^{\infty}I_n,\qquad F=(I_n)_{n\in\mathbb{N}},\qquad
P_{\mathbb{R}}=\mathsf{OrderedSet}(\mathbb{R},\leq),
\]
\[
C=\varnothing
\Longrightarrow
\neg\operatorname{NestedClosedBoundedIntervals}(F,P_{\mathbb{R}}).
\]
\end{remark*}

\begin{remark*}[Failure modes]
\begin{description}
  \item[Exposition.]
  Failure is witnessed by the displayed negated or contrapositive condition.

  \item[\operatorname{MembersSatisfy} / \operatorname{Nested}.]
  \[
\forall I_n\;\forall\mathcal{I}\;
\left[
\mathcal{I}=\mathsf{IndexedFamily}(I_n,\mathbb{N},\mathbb{R})
\land
\operatorname{MembersSatisfy}(\mathcal{I},\mathsf{NonemptyBoundedInterval},\mathbb{R})
\land
\operatorname{Nested}(\mathcal{I},\subseteq)
\Longrightarrow
\bigcap_{n=1}^{\infty}I_n\ne\varnothing
\right].
\]
\[
C=\bigcap_{n=1}^{\infty}I_n,\qquad F=(I_n)_{n\in\mathbb{N}},\qquad
P_{\mathbb{R}}=\mathsf{OrderedSet}(\mathbb{R},\leq),
\]
\end{description}
\end{remark*}

\begin{remark*}[Interpretation]
Dropping closedness can remove the endpoint that would otherwise remain in
the limiting intersection.
\end{remark*}

\begin{dependencies}
\hyperref[def:open-interval]{Open Interval};
\hyperref[def:bounded-subset-of-r]{Bounded Subset of the Real Line}.
\end{dependencies}

\begin{propositionbox}{Proposition (Boundedness Is Necessary in the Nested Interval Theorem)}
\begin{proposition}[Boundedness Is Necessary in the Nested Interval Theorem]
\label{prop:boundedness-is-necessary-in-nested-interval-theorem}
There exists a sequence $(I_n)_{n\in\mathbb{N}}$ of nonempty closed unbounded
intervals in $\mathbb{R}$ such that
\[
  (\forall n\in\mathbb{N})(I_{n+1}\subseteq I_n),
\]
but
\[
  \bigcap_{n=1}^{\infty}I_n=\varnothing.
\]
For example,
\[
  I_n=[n,\infty).
\]
\hyperref[prf:boundedness-is-necessary-in-nested-interval-theorem]{Go to proof.}
\end{proposition}
\end{propositionbox}

\begin{remark*}[Standard quantified statement]
\[
\exists I_n\;
\bigl(
\operatorname{NestedClosedIntervals}(I_n)
\land
\neg\operatorname{IsBoundedInterval}(I_n)
\land
\bigcap_{n=1}^{\infty}I_n=\varnothing
\bigr).
\]
\end{remark*}

\begin{remark*}[Predicate reading]
\[
F=(I_n)_{n\in\mathbb{N}},\qquad C=\bigcap_{n=1}^{\infty}I_n,\qquad
\neg\operatorname{IsBoundedInterval}(I_n,\mathbb{R})
\quad\text{can make}\quad
\operatorname{IsNonempty}(C)
\text{ fail.}
\]
\end{remark*}

\begin{remark*}[Negated quantified statement]
\[
\forall I_n\;
\left[
\operatorname{NestedClosedIntervals}(I_n)
\Longrightarrow
\bigcap_{n=1}^{\infty}I_n\ne\varnothing
\right].
\]
\end{remark*}

\begin{remark*}[Negation predicate reading]
\[
C=\bigcap_{n=1}^{\infty}I_n,\qquad F=(I_n)_{n\in\mathbb{N}},\qquad
P_{\mathbb{R}}=\mathsf{OrderedSet}(\mathbb{R},\leq),
\]
\[
C=\varnothing
\Longrightarrow
\neg\operatorname{NestedClosedBoundedIntervals}(F,P_{\mathbb{R}}).
\]
\end{remark*}

\begin{remark*}[Failure modes]
\begin{description}
  \item[Exposition.]
  Failure is witnessed by the displayed negated or contrapositive condition.

  \item[\operatorname{NestedClosedIntervals}.]
  \[
\forall I_n\;
\left[
\operatorname{NestedClosedIntervals}(I_n)
\Longrightarrow
\bigcap_{n=1}^{\infty}I_n\ne\varnothing
\right].
\]
\[
C=\bigcap_{n=1}^{\infty}I_n,\qquad F=(I_n)_{n\in\mathbb{N}},\qquad
P_{\mathbb{R}}=\mathsf{OrderedSet}(\mathbb{R},\leq),
\]
\end{description}
\end{remark*}

\begin{remark*}[Interpretation]
Nested closed bounded intervals cannot keep shrinking or relocating in a way
that loses every real point. The endpoint lemmas convert containment into
monotone endpoint control, and completeness supplies a real point that survives
all stages at once.
\end{remark*}

\begin{dependencies}
\hyperref[def:closed-interval]{Closed Interval}.
\end{dependencies}

\begin{figure}[H]
\centering
\begin{tikzpicture}[x=0.95cm, y=0.9cm, >=Latex]
  % Draw four nested intervals, each row shifted down
  \foreach \row/\la/\rb/\label in {
    0/0.0/10.0/{$[a_1, b_1]$},
    1/1.2/8.5/{$[a_2, b_2]$},
    2/2.0/7.0/{$[a_3, b_3]$},
    3/2.6/5.8/{$[a_4, b_4]$}
  }{
    \draw[thick, blue!70!black] (\la, -\row*0.8) -- (\rb, -\row*0.8);
    \fill[blue!70!black] (\la, -\row*0.8) circle (2.0pt);
    \fill[blue!70!black] (\rb, -\row*0.8) circle (2.0pt);
    \node[left] at (\la-0.1, -\row*0.8) {\small \label};
  }
  % Indicate left endpoints climbing
  \draw[->, dashed, red!70!black, thick] (0.0,-0.0) -- (0.0,-3.2);
  \node[right, red!70!black] at (0.05, -1.6) {\small $(a_n)\!\uparrow$};
  % Indicate right endpoints falling
  \draw[->, dashed, red!70!black, thick] (10.0,0.0) -- (10.0,-3.2);
  \node[left, red!70!black] at (9.95, -1.6) {\small $(b_n)\!\downarrow$};
  % Pinch point
  \draw[dotted] (4.0, 0.4) -- (4.0, -3.3);
  \node[above] at (4.0, 0.45) {\small $x \in \bigcap [a_n,b_n]$};
\end{tikzpicture}
\caption{Nested intervals $[a_1,b_1] \supseteq [a_2,b_2] \supseteq \cdots$ pinch toward a
common point. The sequence of left endpoints is increasing and bounded above; its
supremum $x$ lies in every interval.}
\label{fig:nested-intervals}
\end{figure}
